\section{Related Work}
\begin{table}[t]
\centering
\caption{Comparison of EHR benchmarks relevant to clinical reasoning and agentic evaluation.
\textbf{Code Gen.}: requires generating executable code or SQL.
\textbf{Scale} ($>$10K): $>$10{,}000 QA instances.
\textbf{Agentic}: supports iterative agent loop with execution feedback.
\textbf{Compositional}: explicitly tests multi-hop reasoning over dependent clinical concepts.
\textbf{Longitudinal}: evaluates reasoning over patient state as it evolves over time.}
\label{tab:benchmark_comparison}
\resizebox{\textwidth}{!}{%
\begin{tabular}{lccccc}
\toprule
\textbf{Benchmark} &
\textbf{Code Gen.} &
\textbf{Scale ($>$10K)} &
\textbf{Agentic} &
\textbf{Compositional} &
\textbf{Longitudinal} \\
\midrule
MIMICSQL~\cite{wang2020treqs}          & \cmark & \cmark & \xmark & \xmark & \xmark \\
EHRSQL~\cite{lee2022ehrsql}            & \cmark & \xmark & \xmark & \xmark & \xmark \\
DrugEHRQA~\cite{bardhan-etal-2022-drugehrqa} & \cmark & \cmark & \xmark & \xmark & \xmark \\
EHRXQA~\cite{bae2023ehrxqa}            & \cmark & \xmark & \xmark & \xmark & \xmark \\
EHR-SeqSQL~\cite{bae2024ehrseqsql}    & \cmark & \xmark & \cmark & \xmark & \xmark \\
EHRNoteQA~\cite{kim2024ehrnoteqa}     & \xmark & \xmark & \xmark & \xmark & \xmark \\
EHRAgent~\cite{shi2024ehragent}       & \cmark & \xmark & \cmark & \xmark & \xmark \\
MedAgentGym~\cite{xu2025medagentgym}  & \cmark & \cmark & \cmark & \xmark & \xmark \\
MedAgentBench~\cite{kim2025medagentbench} & \xmark & \xmark & \cmark & \xmark & \xmark \\
AgentClinic~\cite{schmidgall2024agentclinic} & \xmark & \xmark & \cmark & \xmark & \xmark \\
EHRStruct~\cite{zhang2025ehrstruct}   & \cmark & \xmark & \xmark & \xmark & \xmark \\
AgentEHR~\cite{liao2026agentehr}      & \xmark & \cmark & \cmark & \xmark & \xmark \\
EHRSHOT~\cite{wornow2023ehrshot}      & \xmark & \cmark & \xmark & \xmark & \xmark \\
\midrule
\textbf{\ours{} (ours)} & \cmark & \cmark & \cmark & \cmark & \cmark \\
\bottomrule
\end{tabular}%
}
\end{table}
\paragraph{EHR Databases and Benchmarks}
The MIMIC database family is the standard resource for EHR research~\cite{johnson2020mimic, johnson2018mimic}. Early text-to-SQL benchmarks over MIMIC-III such as MIMICSQL~\cite{wang2020treqs} and EHRSQL~\cite{lee2022ehrsql} rely on templated or clinician-collected SQL and cover a constrained query space. EHR-SeqSQL~\cite{bae2024ehrseqsql} extends this to multi-turn SQL interactions. Beyond SQL, DrugEHRQA~\cite{bardhan-etal-2022-drugehrqa} provides 70K medication-focused QA pairs over MIMIC-III spanning both structured tables and free-text notes, while EHRXQA~\cite{bae2023ehrxqa} introduces multimodal QA over MIMIC-IV by coupling structured records with chest X-ray images. EHRNoteQA~\cite{kim2024ehrnoteqa} targets free-text discharge summaries, requiring synthesis across multiple clinical notes. EHRSHOT~\cite{wornow2023ehrshot} provides 15 few-shot clinical prediction tasks over longitudinal Stanford EHR data, evaluating foundation models on structured time-series inputs, though tasks are point-in-time predictions rather than agent reasoning. On the code-generation side, EHRAgent~\cite{shi2024ehragent} frames complex tabular EHR reasoning as iterative code generation with execution feedback, and EHRStruct~\cite{zhang2025ehrstruct} benchmarks 11 structured-data tasks spanning retrieval, aggregation, and temporal reasoning over Synthea and eICU. Across these works, evaluation is conducted over static, individually posed queries; none provides mechanisms to assess multi-hop concept dependencies or time-evolving clinical state. Furthermore, all prior benchmarks that involve tool use rely on toolboxes pre-designed by domain experts, leaving open the question of whether LLMs can autonomously generate the libraries needed to answer such questions. Table~\ref{tab:benchmark_comparison} summarizes how these benchmarks compare. \ours{} is the only benchmark to jointly support compositional difficulty stratification (via explicit dependency-chain depth across 63 clinical concepts) and longitudinal reasoning over patient trajectories.

\paragraph{LLM-Based Medical Agents}
Several agent systems interact with EHR data to support clinical decisions~\cite{liao2026agentehr, tang2024medagents, hager2024evaluation, bani2025language}. MedAgentGym~\cite{xu2025medagentgym} provides a code-based gym for training medical agents but requires a separately curated environment per task. Biomni~\cite{huang2025biomni} discovers biomedical tools from literature but relies on large proprietary LLMs with no training. MedAgentBench~\cite{kim2025medagentbench} evaluates agents against a realistic FHIR-based virtual EHR environment with 300 tasks spanning ten clinical categories, while AgentClinic~\cite{schmidgall2024agentclinic} tests multi-turn diagnostic reasoning through simulated patient interaction. A key limitation shared across these systems is the reliance on user-designed tool environments: the clinical functions agents invoke are hand-crafted by researchers, making evaluation contingent on the completeness and correctness of those tools. This design also incurs significant token overhead when agents must re-execute full tool chains for each new query. Broad evaluations such as MedHELM~\cite{bedi2025medhelm, bedi2026holistic} assess static task performance without agentic EHR interaction. Our work is orthogonal. Rather than evaluating an agent's behavior within a fixed tool environment, we ask how well LLMs can generate the tool environment itself.

\paragraph{Autoformalization and Code-Augmented Reasoning}
Autoformalization, the translation of informal statements into formal, executable representations, was demonstrated for mathematics by Wu et al.~\cite{wu2022autoformalization} using LLMs to translate theorems into Lean and Isabelle. We adapt this concept to clinical guidelines, replacing proof checkers with QA-based verification and targeting a Python function library rather than formal proofs. This reframing is key: unlike mathematical theorems, clinical guidelines encode institutional norms that vary across hospitals, so the formalization process itself resolves the definitional ambiguity that makes zero-shot evaluation unreliable. On the reasoning side, ReACT~\cite{yao2023react} showed that interleaving reasoning with tool calls improves agent performance, and program-aided approaches such as PAL~\cite{gao2023pal} and Codex~\cite{chen2021codex} demonstrate the value of code as a reasoning medium. Our approach uses a ReACT-style loop \emph{offline} to generate a stable, reusable function library, separating the costly library-construction phase from inference and avoiding the repeated token expenditure of agents that re-derive the same computations at query time.
